\appendix

\chapter{Source Code Availability}

The complete source code for TechPulse is maintained on GitHub:

\begin{center}
\url{https://github.com/EZ2-555555/data-pipeline}
\end{center}

\noindent A live deployment is accessible at:

\begin{center}
\url{http://techpulse-dev-frontend-939514668437.s3-website-us-east-1.amazonaws.com/}
\end{center}

\noindent A screen-recording walkthrough of the live system is available at:

\begin{center}
\url{https://youtu.be/iQ6LFZnk8F8}
\end{center}

The repository includes: five data ingesters (ArXiv, Hacker News, DEV.to, GitHub Trending, RSS); 17-step preprocessing with tiktoken chunking; fastembed (ONNX Runtime) embedding with \texttt{sentence-transformers/all-MiniLM-L6-v2}; hybrid retrieval (metadata filter, pgvector, BM25, weighted RRF, cross-encoder reranking); RAG orchestration with three-layer hallucination verification; four-backend LLM fallback chain; dual-criteria drift detection; AWS SAM infrastructure-as-code; React frontend; paired evaluation framework (50 queries $\times$ 2 configurations) with RAGAS and statistical tests; four-stage CI/CD via GitHub Actions; 208 unit tests across 15 modules; and Docker Compose local development (6 services).

\chapter{System Setup and Reproduction Steps}
\label{app:setup}

This appendix provides step-by-step instructions for reproducing the TechPulse system in both local development and AWS cloud environments.

\section*{Prerequisites}

\begin{itemize}[leftmargin=0.5in]
  \item Python 3.11+ and \texttt{pip}
  \item Docker and Docker Compose
  \item Node.js 18+ and npm (for frontend)
  \item AWS CLI v2 and AWS SAM CLI (for cloud deployment)
  \item An AWS account (Free Tier eligible)
  \item A Groq API key (free at \url{https://console.groq.com})
\end{itemize}

\section*{Step 1: Clone and Install Dependencies}

\begin{lstlisting}[language=bash, basicstyle=\ttfamily\small]
git clone https://github.com/EZ2-555555/data-pipeline.git
cd data-pipeline
python -m venv .venv
# On Windows: .venv\Scripts\activate
# On macOS/Linux: source .venv/bin/activate
pip install -r requirements.txt
\end{lstlisting}

\section*{Step 2: Local Development with Docker Compose}

Docker Compose starts six services: PostgreSQL with pgvector, pgAdmin, LocalStack (S3 + SQS), FastAPI backend, React frontend, and the ingestion scheduler.

\begin{lstlisting}[language=bash, basicstyle=\ttfamily\small]
# Create a .env file with required variables
cat > .env <<EOF
GROQ_API_KEY=your_groq_api_key_here
LLM_BACKEND=groq
EOF

# Start all services
docker compose up --build -d

# Verify services are running
docker compose ps
\end{lstlisting}

\noindent The frontend is accessible at \texttt{http://localhost:3000} and the API at \texttt{http://localhost:8000}. The database admin interface (pgAdmin) is at \texttt{http://localhost:5050} (login: \texttt{admin@techpulse.dev} / \texttt{admin}).

\section*{Step 3: Initialize the Database}

\begin{lstlisting}[language=bash, basicstyle=\ttfamily\small]
# Initialize schema (creates documents, chunks,
# drift_baselines tables and pgvector HNSW index)
python -c "from src.db.init_schema import init_schema; \
           init_schema()"
\end{lstlisting}

\section*{Step 4: Run Ingestion and Build Corpus}

\begin{lstlisting}[language=bash, basicstyle=\ttfamily\small]
# Run a single ingestion cycle for all five sources
python -c "from src.orchestrator.orchestrator import \
           run_orchestrator; run_orchestrator()"

# Or run the scheduler for continuous ingestion
python -m src.scheduler
\end{lstlisting}

\section*{Step 5: Run Tests}

\begin{lstlisting}[language=bash, basicstyle=\ttfamily\small]
pip install -r requirements-test.txt
pytest tests/ -v --tb=short
# Expected: 208 tests passed across 15 modules
\end{lstlisting}

\section*{Step 6: Run Evaluation}

\begin{lstlisting}[language=bash, basicstyle=\ttfamily\small]
# Execute the 100-sample evaluation
# (50 baseline + 50 hybrid queries)
python evaluation/run_eval.py

# Results are written to:
#   evaluation/results/eval_summary.json
#   evaluation/results/ragas_scores.json
#   evaluation/results/raw_results.json
\end{lstlisting}

\section*{Step 7: AWS Cloud Deployment}

\begin{lstlisting}[language=bash, basicstyle=\ttfamily\small]
# 1. Configure AWS credentials
aws configure

# 2. Edit SAM configuration with your VPC/subnet IDs
#    and database password
vi infra/samconfig-freetier.toml

# 3. Validate configuration (checks for placeholders)
make validate-config

# 4. Build and push Lambda container image to ECR
#    (see infra/README.md for detailed ECR steps)
docker build -f Dockerfile.lambda -t techpulse-lambda .

# 5. Deploy with SAM
sam build --template-file infra/template-freetier.yaml
sam deploy --config-file infra/samconfig-freetier.toml

# 6. Deploy frontend to S3
./deploy-frontend.sh
\end{lstlisting}

\noindent The SAM template provisions all required resources: RDS PostgreSQL (\texttt{db.t3.micro}), S3 buckets, SQS queue with DLQ, four Lambda functions, HTTP API Gateway, SNS alert topic, EventBridge schedules, and CloudWatch alarms. The CI/CD pipeline (GitHub Actions) automates this deployment on merge to \texttt{main}.

